% LOI TUA --- Quyển XXVI
\chapter*{Lời Tựa}
\addcontentsline{toc}{chapter}{Lời Tựa}
\markboth{Lời Tựa}{Lời Tựa}

\begin{truyen}{Lời Tựa}{Parents and Children: Learning to Speak Again}
\chuhoa{M}{ối quan hệ giữa cha mẹ và con cái là một trong những mối quan hệ quan trọng nhất đời người, nhưng cũng là nơi nhiều hiểu lầm lặng lẽ tích lại nhất. Càng thương nhau, nhiều người lại càng dễ làm đau nhau bằng lời nói vội, kỳ vọng nặng, hoặc im lặng quá lâu.}
\textit{(The relationship between parents and children is one of the most important in life, but also one of the places where misunderstandings quietly pile up. The more people love each other, the more easily they may hurt each other through rushed words, heavy expectations, or long silence.)}

Quyển hai mươi sáu chọn chủ đề \textbf{phụ huynh và con cái} như một phần nối dài tự nhiên từ những câu chuyện lớp học. Bởi phía sau một học sinh trong lớp luôn là một gia đình ở nhà, với yêu thương, lo lắng, áp lực và những câu chưa nói hết.
\textit{(Volume twenty-six takes on the theme of \textbf{parents and children} as a natural continuation of classroom stories. Behind every student in class there is always a family at home, carrying love, worry, pressure, and many unfinished sentences.)}

Trong 10 chương và 100 màn đối thoại, bạn sẽ gặp những cuộc nói chuyện về việc học, sự so sánh, tự do, kỷ luật, tiền bạc, tổn thương, trưởng thành, và những câu hỏi không có đáp án hoàn hảo. Điều quan trọng không phải là thắng nhau trong lời nói, mà là hiểu nhau thêm một chút sau mỗi lần đối thoại.
\textit{(Across 10 chapters and 100 dialogues, you will meet conversations about study, comparison, freedom, discipline, money, hurt, growing up, and questions without perfect answers. What matters is not winning a conversation, but understanding each other a little more afterward.)}
\end{truyen}

\begin{baihoc}
Nhiều gia đình không cần thêm tiền hay thêm lý lẽ trước tiên. Họ cần thêm khả năng nói với nhau cho đúng.
\textit{(Many families do not first need more money or more arguments. They need a better way to speak to each other.)}
\end{baihoc}

\begin{ghinhoanh}
\textbf{Short English to remember:}

Love is not always the problem.

Sometimes the problem is the way love is spoken.
\end{ghinhoanh}

\begin{tuhocnhanh}
1. Trong nhà em, điều gì thường được thương nhưng ít được nói đúng? \textit{(What is often loved but poorly expressed in your home?)}

2. Em muốn cha mẹ hiểu mình ở điểm nào nhất? \textit{(Where do you most want your parents to understand you?)}

3. Nếu là phụ huynh, em nghĩ mình sẽ cần học lại điều gì trong cách nói chuyện với con? \textit{(If you were a parent, what would you need to relearn in speaking with a child?)}
\end{tuhocnhanh}

\begin{flushright}
\textbf{Tôi Nguyễn Văn Sang}\\
\textit{GV THPT Nguyễn Hữu Cảnh - TP HCM}
\end{flushright}

\ngancach